\section{Methodology}

Driven by the observations, we propose \dawn{}, a training-free, dependency-aware solution to accelerate dLLM inference while maintaining the generation quality.  
% The key idea is to extract positional dependencies to steer dLLMs along more efficient decoding trajectories that avoid errors arising from nonindependent token probabilities across positions. 
The key idea is to \textit{extract positional dependencies and leverage them to select a larger set of reliable positions for update at each iteration}.
\dawn{} realizes efficient inference by three cooperating modules, as shown in Fig~\ref{fig:arch}.
\textit{Dependency Graph Construction} builds a sparse directed graph from a lightweight proxy of token dependencies based on attention maps. \textit{Anchor-Guided Decoding} first selects high-confidence masked positions that are approximately independent, and then leverages high-confidence prompts or committed positions as anchors, enabling strongly coupled positions (induced positions) to be unmasked with a relatively low confidence. These together produce a parallel set $\mathcal{U}^{(t)}_{\text{anchor}}$. For the remaining tokens, \textit{Conflict-Based Scheduling} uses conflict relations induced by the dependency graph to select a maximum independent set $\mathcal{U}^{(t)}_{\text{conflict}}$ from positions that meet a lower confidence threshold. Finally, the positions in  $\mathcal{U}^{(t)}_{\text{anchor}}\cup~\mathcal{U}^{(t)}_{\text{conflict}}$ are unmasked in parallel in this iteration. These three modules are pipelined within each denoising iteration to enable accurate and efficient inference. The following subsections describe their mechanisms in detail.

\subsection{Dependency Graph Construction}
Many prior works \cite{fastv, spargeattn, sageattention, svg} use attention maps to characterize interactions among tokens, thereby enabling more efficient inference. Inspired by these applications, we treat attention maps as a lightweight, approximate proxy for token dependencies during decoding. This coupling signal is readily available from each forward pass and can be converted into a sparse directed dependency graph, which serves as the basis for subsequent efficient scheduling.

Since attention patterns evolve across denoising iterations, the dependency graph is constructed at each iteration from the model's attention maps. To obtain a signal that is closer to the final prediction and less noisy, attention weights are averaged across the last few layers and all heads. As discussed in Sec.~\ref{obs: attnsinks}, dLLMs often exhibit attention sinks, where a small set of positions absorbs most of the attention mass due to systematic bias rather than semantic relevance. This can lead to misleading dependent relations when attention is used as a proxy for dependency. To mitigate this effect, sink positions are identified via outlier detection and filtering of their incoming attention mass: given an aggregated attention matrix $A^{(t)}\in\mathbb{R}^{L\times L}$ at iteration $t$, the incoming attention mass of position $j$ is defined as
\begin{equation*}
\bar{A}^{(t)}_j = \frac{1}{L}\sum_{i=1}^{L} A^{(t)}_{i,j},
\end{equation*}
and position $j$ is marked as a sink if $\bar{A}^{(t)}_j$ is larger than a predefined threshold $\tau_{sink}$. Self-attention on the diagonal is ignored as it does not capture cross-position dependencies.

Given the processed attention proxy at iteration $t$, a directed sparse dependency graph is constructed to capture salient token couplings during decoding. 
Specifically, to keep the graph sparse and focus on the most informative relations, edges are retained based on thresholded ($\tau_{edge}$) attention scores. A directed edge $j\!\rightarrow\! i$ is added if query position $i$ places a sufficiently large attention mass on key position $j$, indicating that the prediction at token $i$ can be significantly conditioned by token $j$. 
 The resulting graph provides an approximate representation of positional dependencies and will be used by subsequent modules.

\subsection{Anchor-Guided Decoding}
Prior work \cite{fastdllm} suggests that sufficiently high-confidence positions are approximately independent as well as non-conflicting, motivating a high conservative threshold $\tau_{high}$ for selecting them for parallel updates, where $\tau_{high}$ is usually set to 0.9. 
Meanwhile, some low-confidence tokens at masked positions can still be consistent with the final response, indicating that only confidence thresholding remains a limit for dLLM inference. Motivated by this, \textit{Anchor-Guided Decoding} is introduced to accept approximately independent positions and reliable positions under a relaxed confidence criterion. 

%After selecting positions whose confidence passes $\tau_{high}$. 
As discussed in Sec.~\ref{obs: anchors}, low-confidence positions that are strongly coupled to previously unmasked positions meet the confidence threshold $\tau_{high}$ and often match the final result. Following this observation, we define anchors as positions that have been unmasked and their confidence meets $\tau_{high}$. Using the dependency graph, we identify the induced positions $\mathcal{I}$ as masked positions that are reachable via directed edges from anchor tokens. Intuitively, higher-confidence anchors provide more reliable context, allowing the corresponding induced positions to be unmasked at lower confidence. Accordingly, the set of positions eligible to be unmasked at the denoising step $t$ is defined as:
% \begin{equation*}
% \mathcal{U}^{(t)}_{\text{anchor}} = \left\{\, i \in M^{(t)} \; \middle| \;ci \ge \tau_{high}\,\right\} \cup
% \left\{\, i \in \mathcal{I}^{(t)} \; \middle| \;
% c_i \ge \tau_{induced} \,\right\},
% \end{equation*}
% \begin{equation*}
% \mathcal{U}^{(t)}_{\text{anchor}}
% =
% \left\{
% \, i \;\middle|\;
% \big(i \in M^{(t)} \land c_i \ge \tau_{\text{high}}\big)
% \ \lor\
% \big(i \in \mathcal{I}^{(t)} \land c_i \ge \tau_{\text{induced}}\big)
% \right\}.
% \end{equation*}
\begin{equation*}
\mathcal{U}^{(t)}_{\text{anchor}}
=
\left\{
\, i \;\middle|\;
\begin{array}{l}
i \in M^{(t)},\ c_i \ge \tau_{\text{high}}
\ \text{or}\\
i \in \mathcal{I}^{(t)},\ c_i \ge \tau_{\text{induced}}
\end{array}
\right\}.
\end{equation*}
where $\tau_{induced}$ is the confidence threshold for unmasking induced positions. 
Overall, \textit{Anchor-Guided Decoding} selects approximately independent positions under a high confidence threshold and relaxes the threshold for induced positions, enabling more efficient inference.
% Overall, \textit{Anchor-Guided Decoding} enables more efficient unmasking while maintaining decoding quality by allowing certain reliable positions to be unmasked even when their confidence falls below the base unmasking threshold.

\subsection{Conflict-Based Scheduling}
% Prior work \cite{fastdllm} suggests that sufficiently high-confidence positions are approximately independent as well as non-conflicting, motivating a conservative threshold for selecting them for parallel updates.
For the remaining positions that are strongly coupled and satisfy a lower confidence threshold $\tau_{low}$, simultaneously unmasking them may introduce inconsistent commitments, as discussed in Sec.~\ref{prl: parallel}. To mitigate such errors while retaining high parallelism, \textit{Conflict-Based Scheduling} is introduced to prevent lower-confidence but highly coupled positions from being decoded in the same step.

Since the dependency graph captures salient positional dependencies, it can be used to identify strongly coupled position pairs. A conflict relation is defined between two positions connected by an edge in the dependency graph, regardless of direction: if either $i\rightarrow j$ or $j\rightarrow i$ is present, then positions $i$ and $j$ are considered to be in conflict and should not be unmasked simultaneously.

Algorithm~\ref{alg:conflict} shows the detailed procedure. Specifically, at iteration $t$, based on these conflicts and the graph topology, a greedy independent set is constructed from the remaining positions that (i) are not yet included in and not conflicting neighbors of ~$\mathcal{U}^{(t)}_{\text{anchor}}$, and (ii) satisfy a lower confidence threshold, yielding additional positions that can be decoded in parallel. Concretely, positions are greedily selected in descending order of confidence: the highest-confidence position is added to the update set $\mathcal{U}^{(t)}_{\text{conflict}}$, and all of its conflicting neighbors are removed from further consideration. This procedure repeats until no candidate positions remain.

\begin{algorithm}[tb]
  \caption{Conflict-Based Scheduling}
  \label{alg:conflict}
  \begin{algorithmic}[1]
    \STATE {\bfseries Input:} remaining candidate positions $\mathcal{C}^{(t)}$, anchor update set $\mathcal{U}^{(t)}_{\text{anchor}}$, confidence scores $\{c_i\}_{i\in\mathcal{C}^{(t)}}$, conflict neighbors $\{\mathcal{N}(i)\}_{i\in\mathcal{C}^{(t)}}$, lower threshold $\tau_{\mathrm{low}}$
    \STATE {\bfseries Output:} parallel update set $\mathcal{U}^{(t)}_{\text{conflict}}$
    \STATE Initialize $\mathcal{U}^{(t)}_{\text{conflict}} \leftarrow \emptyset$
    \STATE $\mathcal{X} \leftarrow \mathcal{U}^{(t)}_{anchor} \cup \bigcup_{i \in \mathcal{U}^{(t)}_{anchor}} \mathcal{N}(i)$ \COMMENT{selected positions and their conflicts}
    \STATE $\mathcal{R} \leftarrow \{\, i \in \mathcal{C}^{(t)} \mid c_i \ge \tau_{\mathrm{low}} \,\} \setminus \mathcal{X}$ \COMMENT{remaining candidates}
    \WHILE{$\mathcal{R} \neq \emptyset$}
      \STATE Select $i^\star \leftarrow \arg\max_{i \in \mathcal{R}} c_i$
      \STATE $\mathcal{U}^{(t)}_{\text{conflict}} \leftarrow \mathcal{U}^{(t)}_{\text{conflict}} \cup \{i^\star\}$
      \STATE $\mathcal{R} \leftarrow \mathcal{R} \setminus \left(\{i^\star\} \cup \mathcal{N}(i^\star)\right)$
    \ENDWHILE
    \STATE {\bfseries return} $\mathcal{U}^{(t)}_{\text{conflict}}$
  \end{algorithmic}
\end{algorithm}

In practice, quality degradation from naively lowering the confidence threshold is largely driven by positional coupling. By explicitly avoiding simultaneous unmasking of strongly coupled positions, \textit{Conflict-Based Scheduling} helps maintain decoding quality while allowing a lower confidence threshold $\tau_{low}$, thus speeding up the inference. 


